\section{The weakened step quantities}\label{sec:step-quantities}

This section turns the ingredients of one step (Theorem~\ref{thm:step}) into the input of the cluster expansion. That input consists of functions of the weakening parameters $s\in\mathbb D_\kappa$, the white-noise field and the background, with three properties. They are analytic with uniform bounds; they factorize when $s_\Box=0$ outside a union of boxes; and they are invariant under real gauge transformations. They are defined by source-free formulas for every background in the domain, not only at critical backgrounds. Section~\ref{sec:cluster} integrates them, and the flow of Section~\ref{sec:flow4} applies the result at every step.

The determinant is treated in Proposition~\ref{prop:one-loop-polymer}, and the Gaussian measure through $C_k^{1/2}(s)$ (Corollary~\ref{cor:weak-covariance}). This section treats the interaction (Section~\ref{subsec:sq-interaction}), the gauge-fixing remainder (Section~\ref{subsec:sq-gf}), the Jacobians (Section~\ref{subsec:sq-jacobians}), the old activities composed with the background map (Section~\ref{subsec:sq-composed}) and the Hessian away from critical backgrounds (Section~\ref{subsec:sq-hessian}), which Section~\ref{sec:localization} already defines. Section~\ref{subsec:sq-loc} gives the localization lemma that turns these functions of $s$ into pieces of volume type.

\paragraph{The formula of one step.}\phantomsection\label{par:step-formula}%
Recall the step of Part~VII. The reference $U'=U_{k+1}(V')$ is a background field of the configuration $V'$ of step $k+1$, the base $W=V^0=U'^{(k)}$ is its average at scale $k$ (Lemma~\ref{lem:composition}), and the fluctuation field $Z$ lives on the bonds $\mathfrak F$ of Definition~\ref{def:bonds}; the step integrates over the configurations $e^{\ii\hat Z(Z)}W$ (Proposition~\ref{prop:exact-step}). Theorem~\ref{thm:step} writes the step as
\[
\begin{aligned}
T^{\rm gf}(\chi^{\rm w}\rho)(V')={}&\mathcal N_k\exp\Big(-\frac1{g^2}S(U')-\frac12\sum_{b\in\mathfrak F}\ell_k(b)+\sum_c\log\mathfrak j_c\Big)\\
&\times\int d\mu_{C_k}(Y)\,\chi^{\rm w}\big(\hat Z(gY)\big)\exp\Big(-\frac{(\mathcal V+G_3)(gY)}{g^2}+\mathfrak J(gY)+\mathfrak E\big(e^{\ii\hat Z(gY)}W\big)\Big).
\end{aligned}
\]
Here $\mathcal N_k$ does not depend on the background (Theorem~\ref{thm:step}(a)); $\ell_k(b)$ are the determinant densities of Proposition~\ref{prop:det-densities}; $\mathfrak j_c$ are the Jacobians of the coarse bonds $c$ (Theorem~\ref{thm:step}(b)); $\mu_{C_k}$ is the Gaussian measure with the covariance $C_k=\Delta_k^{-1}$ of Corollary~\ref{cor:covariance}; $\chi^{\rm w}$ is the fluctuation cutoff of Definition~\ref{def:chi-w}, and $\hat Z$ the elimination map of Lemma~\ref{lem:eliminate}; $\mathcal V+G_3$ is the interaction, of third order (Theorem~\ref{thm:step}(d)); $\mathfrak J$ collects the logarithms of the densities $\sigma$ and of the Jacobians (Theorem~\ref{thm:step}(e)); and $\mathfrak E$ is the sum of the old activities, composed with the background map (Theorem~\ref{thm:step}(f)). In this section the determinant becomes the one-loop pieces of Proposition~\ref{prop:one-loop-polymer}, the measure is written with $C_k^{1/2}(s)$, the interaction gives the pieces of Sections~\ref{subsec:sq-interaction} and~\ref{subsec:sq-gf}, $\mathfrak J$ those of Section~\ref{subsec:sq-jacobians}, and $\mathfrak E$ the composed activities of Section~\ref{subsec:sq-composed}.

\paragraph{Setting.}
\begin{itemize}
\item \emph{The reference.} $U'$ is a reference in the sense of Section~\ref{sec:weak-background}: real and $(M,\eps)$-regular with $|J_{U'}|\le\eps$, or a complex reference $U'=e^{\ii\eta A'}U'_r$ of size $\delta_c<\delta_P$ in the sense of Section~\ref{sec:complex-references}, whose real part is $(M,\eps)$-regular with $\eps\le\eps_{\max}$ (Proposition~\ref{prop:complex-ext}) and satisfies $|J_{U'_r}|\le\eps$. Theorem~\ref{thm:weak-bg}(a), (b), (d) hold at it (Proposition~\ref{prop:complex-ext}(d)). The norm $|\cdot|_*$ is that of Theorem~\ref{thm:weak-bg}, with $\beta=\alpha$, and the H\"older exponent $\alpha$ used below is fixed with $\alpha\le\frac12$, as in Section~\ref{sec:weak-background}. The constant $K_H$ and the radii $R_P$, $r_P$ of Part~IX are those computed with the multiplicative constants doubled (Proposition~\ref{prop:complex-ext}, end of the proof); they serve at every complex reference with $\delta_c<\delta_P$ and do not depend on $\delta_P$. The thresholds $M_1$ and $M'_1$ are those of Remark~\ref{rem:parameters}(4) and do not depend on $\delta_P$. At a complex reference, $|\cdot|_*$ is taken in the real gauges of $U'_r$, and covariant derivatives with respect to $U'_r$; the other reading, with respect to $U'$, differs by a local term of size $C\delta_c|A|$. Within this setting, where regularity is assumed separately, the current bound on $U'_r$ is used at complex references only for coercivity (Proposition~\ref{prop:hessian-general}(d)); at real references it is also used in Theorem~\ref{thm:weak-bg}(c). On the graded domains it is used in addition for the $(M,\eps)$-regularity of $U'_r$ (Corollary~\ref{cor:regularity}), for the background-activity domain (Proposition~\ref{prop:one-loop-polymer}(b)) and in Lemma~\ref{lem:quadratic-bound}. It holds there by (DC2): in the units of step $k$, $|J_{U'_r}|<L^{-3}\theta_r\le\eps$ (Remark~\ref{rem:domain-constants}(ii)).
\item \emph{Extension.} A point of the activity domain of step $k+1$ at a polymer $Y'$ of step $k+1$ fixes $U'$ only near $Y'$. It is extended to a reference by the cutoff of Remark~\ref{rem:bg-to-graded}, with the same real part $U'_r$. A piece indexed by a polymer $Z$ is used only when $Y'$ contains the $LM$-cubes meeting $Z^{(r')}$ (Remark~\ref{rem:extension}). By the factorization below, such a piece does not depend on the extension.
\item \emph{The base.} For every reference put
\[
W:=U'^{(k)},\qquad V':=U'^{(k+1)} .
\]
At critical backgrounds these are $V^0$ and $V'$ of Lemma~\ref{lem:composition}. They are local in $U'$ and covariant, and analytic at complex references by Proposition~\ref{prop:AB}(i). The elimination of Lemma~\ref{lem:eliminate}, for the plain constraint $\ell_c$ of Section~\ref{sec:fluctuations}, is taken at the base $W$. In that lemma the parameter $\mathsf a$ is defined through the curvature $\eps'_1$ of $V'$; at a general reference it is replaced by a bound on the curvatures of the real parts $U'_r$ and $(U'_r)^{(k)}$, which is all the proof uses. On the graded domains these curvatures are at most $2\theta_r$ (Remark~\ref{rem:domain-constants}(ii)), so $L^d\mathsf a\le c$ is (DC3). Since $c\le r_M/64$ in Lemma~\ref{lem:eliminate}, the real part $W_r=(U'_r)^{(k)}$ of the base then satisfies $|\ell(W_r(\partial p))|\le L^{-2}\mathsf a\le\frac18\eps_M$, so it lies in the regular set $\mathfrak M$ of Section~\ref{sec:fluctuations} with the margin of \eqref{eq:base-regular}. At complex bases the elimination is analytic in the reference, by the argument in the proof of Corollary~\ref{cor:covariance} together with Proposition~\ref{prop:AB}(i),(vi). This needs $C_{\rm AB}\delta_c\le c'_L$, which is part of the smallness of $\delta_P$ (Proposition~\ref{prop:complex-ext}). It gives $\zeta_c$, $\hat Z$, $\hat Z^1$, $\hat Z^2$, with $|\hat Z(Z)|_\infty\le CL^d|Z|_\infty$. The averaged contour variables of the step and the transports $\mathfrak G_W$ at complex bases are defined and holomorphic by (DC5) (Remark~\ref{rem:domain-constants}, reasons for (ii)).
\item \emph{The fields.} We write $\xi$ for the white-noise field on $\mathfrak F$ (Dimock's $W$), since in Sections~\ref{sec:step-quantities}--\ref{sec:target-flow} $W$ is the base, and $Z(s)=gC_k^{1/2}(s)\xi$. $A_*(s;B)$ is the weakened background map of Theorem~\ref{thm:weak-bg} at the reference $U'$.
\item \emph{Margins.} Let $r$ be the largest margin of the weakened families used: $\bC(s)$, $\bH(s)$, $D\bH(s)$ and their transposes, and $C_k^{1/2}(s)$ (Corollary~\ref{cor:weak-covariance}). The definition of weakenable families (Section~\ref{sec:weakenable}) requires $r\le M/16$. The margins are sums of a bounded number of ranges of local operators (Lemma~\ref{lem:weak-local-ops}(a), Theorem~\ref{thm:weak-inverse}(a), and the box split of the additive local terms, Section~\ref{sec:weakenable}), independent of $M$; so $r\le M/16$ is a condition on $M$, listed in Remark~\ref{rem:parameters}. A family with margin $r$ is also a family with every larger margin up to $M/16$, since its kernel supports and its data dependence lie in every larger neighbourhood; so we may and do assume $r\ge3$. Put $r'=r+2L+2$. We assume $M>32(L+1)$, so that $2r'<M/4$. Two polymers that do not touch are then at distance at least $M/4>2r'$, so their $r'$-neighbourhoods are disjoint.
\item \emph{Constants.} $C_P$ denotes constants that may depend on $\kappa,\delta,L,M,\rho,\alpha$, but not on $k$, $n$, $\eta$ or $s$.
\end{itemize}

\begin{remark}[Notation in Sections~\ref{sec:step-quantities}--\ref{sec:flow4}]\label{rem:notation-flow}
Some letters are used here with a meaning that differs from earlier parts of the note. We keep them, and list them once.
\begin{itemize}
\item $C_T$ is the bound on $C_k^{1/2}(s)$ of Section~\ref{subsec:sq-interaction}, not the tier constant of Section~\ref{sec:tiers}. $C_{\rm tor}$ is the torus constant of (T2).
\item $C_d$, $A$, $\epsilon_*$ are the constants of Theorem~\ref{thm:mayer}; they are not the constants of Lemma~\ref{lem:weak-chain} or the bootstrap threshold of the earlier parts.
\item $\kappa$ and $\kappa'=\kappa-C\delta$ are the weakening rates of Part~IX; Section~\ref{sec:flow} writes $\kappa'_w$ for $\kappa'$, since there $\kappa'$ is also the reblocked rate of Lemma~\ref{lem:reblocking}. $\kappa_n$ is the rate of the activity norm, fixed in Remark~\ref{rem:parameters}(1). The rate $\kappa$ in Theorem~\ref{thm:mayer} is generic. Theorem~\ref{thm:weak-bg} is stated on $\mathbb D_\kappa$ with $\kappa$ the rate of its inputs $\bH(s)$, $\bC(s)$ at level~(L4); that rate is at least $\kappa'$, and the pieces below are formed on $\mathbb D_{\kappa'}$.
\item $\ell_k(b;1)$ is the density at the reference $U'=1$, not at $s=1$.
\item $U'$ is the reference; in Section~\ref{sec:complex-references} it was also the relative factor $e^{\ii\eta A'}$. $W$ is the base; earlier it also denotes potentials and configurations.
\item $r$ is the margin; the resolvent parameter of Corollary~\ref{cor:weak-covariance} is unrelated. $r_E$ is the margin of dependence of the old activities (Section~\ref{sec:cluster}).
\item $\delta$ is the loss parameter of Part~IX, $\delta_c$ the size of a complex reference, and $\theta_r,\delta_1,\delta_0$ (with primes) the budgets of a graded domain. $\bar\delta$ is the cap of Remark~\ref{rem:domain-constants}, and $\theta$ in Section~\ref{sec:tiers} is unrelated; $\delta_0$ of Proposition~\ref{prop:reg-complex} enters only through $\delta_P$.
\item $B_P$, $B_3$, $B_5$ are unrelated to the coarse field $B$; $C_{\rm P}$ (roman) is the Poincar\'e constant of Lemma~\ref{lem:poincare-step}, unrelated to $C_P$. In Part~IX, $\mathsf N_{\rm b}$ is the box count of Section~\ref{sec:weakenable}, and $\mathsf N$, in Proposition~\ref{prop:AB} and earlier, is the matrix size of $G$. $h_j$ is a partition of unity, unrelated to $h_k=p_kg_k$. $\mathsf a$ is the curvature parameter of Lemma~\ref{lem:eliminate}, not the averaging parameter $a$.
\item $\hat Z$ is the elimination map acting on fields; applied to a polymer (Remarks~\ref{rem:localization-margins} and~\ref{rem:reblock-margins}) it denotes the union of the $LM$-cubes meeting $Z^{(r')}$. The polymer $Z$ and the field $Z(s)$ are distinguished by the argument.
\item $p$ is a plaquette in $s_p$ and $v_p$, and $p$, $p_0=p_{0,k}$ bound the white noise; $p_k=p(g_k)$ is the small-field scale. $\zeta$ is the bound on $\hat Z$, unrelated to $\zeta_c$. $\tau$ is the composed-activity ratio of Proposition~\ref{prop:composed}. $\varphi_{y,x}$ are the gauge-fixing pieces of Definition~\ref{def:gf-pieces}, $G_3(\Box;s)$ their sum over a box, and $f(H)=\mathsf N-\operatorname{tr}\cos H$; the composed activities are $f(X;s,Z)$, with a polymer as first argument. $\mathsf F$ (Section~\ref{sec:weak-background}) and $F^s$ denote the same symmetrized curvature.
\item $\epsilon$ is, in Section~\ref{sec:target-flow}, the finest spacing; $\eps_k$ is the vacuum energy per volume, unrelated to the thresholds $\eps_0$, $\eps_1$, $\eps_M$ and to $\eps^0$ of $N_\chi$. $n=n_k$ is the side of $\T^{(k)}$ (see (T0)), and $n_0$ is a threshold, not a value at step $0$. In Sections~\ref{sec:target-flow} and~\ref{sec:flow4}, $h_0=p(g_0)g_0$ is the small-field threshold at step $0$; in Section~\ref{sec:AB-proof}, $h_j$ with $h_0=1$ is a different quantity.
\item In Section~\ref{sec:flow4}, $E'=E'^{(k+1)}[\mathcal E_k,g_k]$ is the raw new generation of a step, $Y'_c$ the polymer of scale $k+1$ that carries the Jacobian term of the bond $c$, and $\mathcal E_k(Y)$ the grouped old activities. $B_{\rm new}$, $B^{(0)}_{\rm new}$ bound the new terms of a step, and $\mathsf K=2\max(B_{\rm new},B^{(0)}_{\rm new})$ bounds the generations; $\mathsf K$ is unrelated to the number $K$ of steps, and $B_{\rm new}$ to $B_1$ of Theorem~\ref{thm:chart} and to $B_3$, $B_5$, $B_P$. $C_{\rm old}$, $C_{\mathfrak e}$, $C_{\rm jac}$, $C_{\rm rb}$ and $C_{\rm new}$ are the constants of Proposition~\ref{prop:old-activities} and Lemma~\ref{lem:raw-new}, and ($\Gamma$1)--($\Gamma$4) are the conditions on $\gamma$ of Remark~\ref{rem:flow-constants}; ($\Gamma$4) and its constant $c_{\rm F}$ are used only in Part~XI.
\item In Sections~\ref{sec:generations}--\ref{sec:flow4}: $\mathfrak S^{\rm grid}_j\subset\mathfrak S_j$ is the group of the symmetries of the grid of scale $j$, and $E^{\rm sym}$ the average of a raw generation over it (Definition~\ref{def:symmetrization}); it is not the group $\mathfrak S_j$ of all lattice symmetries. $\mathcal U^\flat$ denotes the domains with $\delta_1^\flat=\frac14\delta_1$ in place of $\delta_1$; $\Theta^\flat$ is the quantity $\Theta$ of Lemma~\ref{lem:later-normal-form} with $\delta_1^\flat$ in place of $\delta_1$ (Remark~\ref{rem:parameters}(5)). $C_{\rm sym}$ is the constant of (DC6). $C_\beta$, $C_{\rm ex}$ and $C_{\rm ne}$ are the constants of Proposition~\ref{prop:period-iso}(b), Proposition~\ref{prop:extraction} and Lemma~\ref{lem:large-later}. $\cX^{\rm sm}_j(\Box)$ is the set of the small polymers of scale $j$ based in $\Box$, $\Psi^{(j)}_\Box$ their sum minus $\beta_jS_\Box$, and $E^{(j),\rm ne}_Y$ the sum over the other polymers of scale $j$ with closure $Y$ (Definition~\ref{def:old-activities}).
\item In Section~\ref{subsec:flow4-flow}: $\Sigma_k$ is the set of the coupling sequences $\pi=(g_0,\dots,g_k)\in(0,\gamma]^{k+1}$ with $g_{i+1}\le2g_i$, not a sum, and $E^{(j)}[\pi]$, $\beta_j[\pi]$ are the generations and the coefficients of a free sequence $\pi$. $\mathsf B_k$ is the set of the bare couplings $g_0$ for which $g_1,\dots,g_k$ are defined and lie in $(0,\gamma]$ (Proposition~\ref{prop:flow-continuity}, Corollary~\ref{cor:bare-continuity}); it is unrelated to $B_{\rm new}$, $B_1$, $B_3$, $B_5$ and $B_P$. $\mathcal A^\flat(Y;U_r)$ and $U_A$ are the fields of the proof of Proposition~\ref{prop:flow-continuity}. $C_{\rm lat}$, $C_{\rm vac}$, $C'_{\rm vac}$ and $\bar p_5$ are the constants of Theorem~\ref{thm:sf-flow}(d),(f) and its proof, and $I_K$, $c_V$, $c'_V$ those of Corollary~\ref{cor:sf-stability} and its proof; $c_V$ and $c'_V$ are unrelated to the coarse configurations $V$ and to the constant $c'$ of Definition~\ref{def:polymer-classes}.
\end{itemize}
\end{remark}

\subsection{The localization lemma}\label{subsec:sq-loc}

For a polymer $Y$, a polymer $Z\supset Y$ and $t\in[0,1]^{Z\setminus Y}$, let $s(Z,t)$ be the point of $\mathbb D_\kappa$ with $s_\Box=1$ for $\Box\subset Y$, $s_\Box=t_\Box$ for $\Box\subset Z\setminus Y$, and $s_\Box=0$ otherwise.

\begin{lemma}[Localization]\label{lem:localization}
Let $\kappa\ge1$, let $Y$ be a polymer, and let $f$ be a function on $\mathbb D_\kappa$, analytic in the interior and continuous, with $|f|\le F$. Suppose that, for every union of boxes $Z'\supset Y$ with $*$-component $Z'_0\supset Y$, and every $s$ with $s_\Box=0$ for $\Box\not\subset Z'$, $f(s)$ depends on $s$ only through $s_\Box$, $\Box\subset Z'_0$, and on the data only on $(Z'_0)^{(r')}$. Then
\[
f(1)=\sum_{Z\supset Y}f_{\rm loc}(Y,Z),\qquad f_{\rm loc}(Y,Z)=\int_{[0,1]^{Z\setminus Y}}\partial_{s_{Z\setminus Y}}f\big(s(Z,t)\big)\,dt ,
\]
the sum running over the polymers $Z\supset Y$, where:
\begin{enumerate}[label=(\alph*)]
\item $f_{\rm loc}(Y,Z)$ depends on the data only on $Z^{(r')}$;
\item $|f_{\rm loc}(Y,Z)|\le F\,e^{-(\kappa-1)|Z\setminus Y|}$;
\item if $f$ is invariant under real gauge transformations of the data, or analytic in further parameters with $|f|\le F$ uniformly, then so is $f_{\rm loc}(Y,Z)$.
\end{enumerate}
\end{lemma}

\begin{proof}\hasdetail{lem:localization}
Write $f(1)$ as a function of the variables $s_\Box$, $\Box\not\subset Y$, with $s_\Box=1$ on $Y$. Apply $g(1)=g(0)+\int_0^1g'(t)\,dt$ in each of these variables. This gives the sum over sets $S$ of boxes outside $Y$ of $\int_{[0,1]^S}\partial_{s_S}f\,dt$, evaluated at $s_\Box=1$ on $Y$, $s_\Box=t_\Box$ on $S$, and $s_\Box=0$ elsewhere. Let $Z'=Y\cup S$, with $*$-component $Z'_0\supset Y$. By hypothesis, at these points $f$ does not depend on $s_\Box$ for $\Box\subset S\setminus Z'_0$. So the term vanishes unless $S\subset Z'_0$, i.e.\ unless $Z=Y\cup S$ is a polymer. This gives the formula. Part~(a) follows from the hypothesis at $Z'=Z$.

For (b), every point $s(Z,t)$ has $|s_\Box|\le1$. In each differentiated variable, Cauchy's formula on the circle of radius $e^\kappa-1$ around $t_\Box$ stays in $\mathbb D_\kappa$, by continuity up to the boundary. Each variable is differentiated once, so $|\partial_{s_S}f|\le F(e^\kappa-1)^{-|S|}\le Fe^{-(\kappa-1)|S|}$ for $\kappa\ge1$. Part~(c) holds because differentiation in $s$ and integration over $t$ preserve both properties.
\end{proof}

\begin{remark}[Margins]\label{rem:localization-margins}
The pieces are indexed by polymers $Z$ and depend on the data on $Z^{(r')}$, not inside $Z$. In the cluster expansion, polymers that do not touch have disjoint $r'$-neighbourhoods, so ``touching'' is the hard-core condition. Margins should not be converted to $0$ for the composed activities. Replacing $Z$ by $Z^\sharp$ (Section~\ref{sec:weakenable}) turns the decay $e^{-(\kappa-1)|Z\setminus Y|}$ into $e^{-3^{-d}(\kappa-1)|Z^\sharp|+(\kappa-1)|Y|}$. For $Y$ a single box, the factor $e^{(\kappa-1)|Y|}$ is a constant. For the composed activities, $Y$ is the old polymer, and the factor would swamp the old decay. The next-step polymer of a piece is the union $\hat Z$ of the $LM$-cubes meeting $Z^{(r')}$. In the reblocking this costs at most $C|\hat Z|_{LM}$, counted in $LM$-cubes, with $C$ independent of $L$ and $M$, since $r'<M/8$.
\end{remark}

\subsection{The interaction pieces}\label{subsec:sq-interaction}

Write $S(U)=\sum_ps_p(U)$, where $s_p$ is the term of the plaquette $p$ of $\Tf$ in the action (Section~\ref{sec:action}); the subscript is a plaquette, and $s_p$ is unrelated to the weakening parameter $s$. Put $r_I=\min\big(r_L,c_IL^{-d}R_P,\frac1{32}c_0L^{-1}\big)$, with $r_L$ from Lemma~\ref{lem:eliminate}, $R_P$ from Theorem~\ref{thm:weak-bg} (computed with the multiplicative constants doubled, see the Setting), $c_0$ from Lemma~\ref{lem:exp-gauge}(c), and $c_I$ so small that $|\hat Z(Z)|_\infty<R_P$ for $|Z|_\infty<r_I$. (For the constant $c$ of Lemma~\ref{lem:eliminate} small, $r_L\le\frac1{32}c_0L^{-1}$, so the third entry does not change the minimum; the second can be smaller than $r_L$, since $R_P$ depends on $\kappa,\delta,M,\rho$.)

\begin{definition}[Interaction pieces]\label{def:interaction-pieces}
For a plaquette $p$, $s\in\mathbb D_\kappa$ and complex $Z$ on $\mathfrak F$ with $|Z|_\infty<r_I$, put
\[
v_p(s;Z)=s_p\big(e^{\ii\eta A_*(s;\hat Z(Z))}U'\big)-T_2\big[s_p\big(e^{\ii\eta A_*(s;\hat Z(\cdot))}U'\big)\big](Z),
\]
where $T_2$ is the Taylor polynomial of degree $2$ at $Z=0$, for fixed $s$ and $U'$.
\end{definition}

\begin{proposition}[Interaction pieces]\label{prop:interaction-pieces}
\begin{enumerate}[label=(\alph*)]
\item $v_p(s;Z)$ is analytic in $(s,Z)$ for $s$ in the interior of $\mathbb D_\kappa$ and $|Z|_\infty<r_I$, continuous up to the boundary in $s$, and analytic in the reference.
\item $|v_p(s;Z)|\le C_P\,\eta^d\,(|Z|_\infty/r_I)^3$ for $|Z|_\infty\le\frac12r_I$, uniformly in $s\in\mathbb D_\kappa$.
\item \emph{(Identification.)} If $U'$ is real and critical, then $\sum_pv_p(1;Z)=\mathcal V(Z)$ for $|Z|_\infty\le\min(\frac12r_I,r_2)$ with $|\hat Z(Z)|_\infty\le\beta_V$.
\item \emph{(Factorization.)} Let $s_\Box=0$ for $\Box\not\subset Z'$, a union of boxes with $*$-components $Z'_i$. Then $v_p=0$ unless $p$ has a point in $Z'^{(r+1)}$. If $p$ has a point in $(Z'_i)^{(r+1)}$, then $v_p$ depends on $s$ only through $s_\Box$, $\Box\subset Z'_i$, on $Z$ only on $(Z'_i)^{(r+2L)}$, and on $U'$ only on $(Z'_i)^{(r')}$.
\item \emph{(Invariance.)} $v_p$ is invariant under real gauge transformations $u$, acting by $U'\mapsto U'^u$ and $Z(b)\mapsto R(u(b_-))Z(b)$.
\end{enumerate}
\end{proposition}

\begin{proof}\hasdetail{prop:interaction-pieces}
(a) follows from Theorem~\ref{thm:weak-bg}(b), the analyticity of $\hat Z$ and of $s_p$, and Proposition~\ref{prop:complex-ext}(d). The Taylor coefficients at $Z=0$ are analytic in $s$ and in the reference.

(b) Fix $s$ and $Z$, and put $g(t)=v_p(s;tZ)$. It is analytic for $|t|<r_I/|Z|_\infty$ and vanishes to third order at $t=0$. Let $h(t)=s_p(e^{\ii\eta A_*(s;\hat Z(tZ))}U')-s_p(U')$. Write $s_p(U)=\mathsf N-\operatorname{tr}\cos X_p$ with $X_p=\ell(U(\partial p))$; in $d=4$ the action carries no power of $\eta$. Then $|s_p(U)|\le C|X_p|^2$ for $|X_p|\le1$. By Theorem~\ref{thm:weak-bg}(a), $|A_*(s;B)|_*\le C_P|B|_\infty\le C_PR_P=C_P\varrho_P/K_H\le\frac18$, since $4C_P(\eps_P+2\varrho_P)\le1$ in its proof (with $K_H\ge1$). So, in the gauge of $U'_r$, the four bond fields of $p$ at $e^{\ii\eta A_*}U'$ have sup norm at most $\frac18+\eps+C\delta_c\le\frac14$, since $\eps+C\delta_c\le\frac1{20}$ (Proposition~\ref{prop:complex-ext}, smallness). For $k\ge1$, $\eta=L^{-k}\le\frac13$, so $|U(\partial p)-1|\le e^\eta-1\le\frac12$. At $k=0$, where $\eta=1$ and $A_*=B$ (Proposition~\ref{prop:first-step}), $|B|_\infty\le\rho_L$ (Lemma~\ref{lem:eliminate}(c)), so, with $\eps+C\delta_c\le\frac1{20}$ and $\rho_L\le\frac1{20}$ (Proposition~\ref{prop:first-step}), the bond fields have sup norm at most $\frac1{10}$ and $|U(\partial p)-1|\le e^{2/5}-1\le\frac12$. In both cases $|X_p|\le1$ (Lemma~\ref{lem:log-estimates}(a)). By the expansion of Lemma~\ref{lem:curvature}, $X_p=\eta^2\big(F_{U'}(p)+O(|A_*|_*)\big)$, and $|F_{U'}|\le C(\eps+\delta_c)$. So $|s_p|\le C\eta^4=C\eta^d$ at $e^{\ii\eta A_*}U'$ and at $U'$, and $|h(t)|\le C_P\eta^d$ for $|t|<r_I/|Z|_\infty$. No expansion in $A_*$ is made, so the current, which is not bounded uniformly in $\eta$ at complex references, does not enter. By Cauchy's estimate each of the three Taylor terms of $h$ of degree at most $2$ is bounded by $C_P\eta^d$ there. So $g=h-T_2h$ is bounded by $4C_P\eta^d$ there, and vanishes to order three. The maximum principle for $g(t)/t^3$ on the discs $|t|\le\varrho$, $\varrho<r_I/|Z|_\infty$, and the limit $\varrho\to r_I/|Z|_\infty$ give $|g(1)|\le4C_P\eta^d(|Z|_\infty/r_I)^3$.

(c) By Theorem~\ref{thm:weak-bg}(c), $A_*(1;B)=A_*(B)$ for $|B|_\infty\le\min(R_P,\beta_V)$. So, summing over the finitely many plaquettes,
\[
\sum_ps_p\big(e^{\ii\eta A_*(1;\hat Z(Z))}U'\big)=\mathcal S(\hat Z(Z)).
\] By the second display of Lemma~\ref{lem:step-quadratic}, $\mathcal S(\hat Z(Z))=S(U')+\frac12\Delta^0(\hat Z^1Z,\hat Z^1Z)+\ip{\omega_0}{\hat Z^2(Z)}+\mathcal V(Z)$, where $\mathcal V$ vanishes to third order. Hence $\sum_pT_2[\cdots]=S(U')+\frac12\Delta_k(Z,Z)-\frac12\ip{\mathfrak G_WZ}{\mathfrak G_WZ}$, and $\sum_pv_p(1;Z)=\mathcal V(Z)$. The missing quadratic term is that of the gauge fixing, which is treated in Section~\ref{subsec:sq-gf}.

(d) By Theorem~\ref{thm:weak-bg}(d), $A_*(s;B)$ vanishes outside $Z'^{(r)}$. On $(Z'_i)^{(r)}$ it depends on $B$ only on the coarse bonds of $(Z'_i)^{(r)}$, on $U'$ only on $(Z'_i)^{(r+1)}$, and on $s$ only in $Z'_i$. If $p$ has no point in $Z'^{(r+1)}$, then $A_*=0$ on the bonds of $p$. So $s_p(e^{\ii\eta A_*}U')=s_p(U')$ for all $Z$, and $v_p=0$. Otherwise, all bonds of $p$ lie in $(Z'_i)^{(r+2)}$ for exactly one $i$, since the components are at distance at least $M/4$. The map $\hat Z$ has range $2L$ in block units, and the base $W=U'^{(k)}$ is local. This gives the dependence stated.

(e) $A_*(s)$ is covariant (Lemma~\ref{lem:weak-basic}(c) and the construction in Section~\ref{sec:weak-background}), $\hat Z$ is covariant, and $s_p$ is invariant. The Taylor coefficients of an invariant function are invariant.
\end{proof}

Grouping by boxes, put $V(\Box;s)=g^{-2}\sum_{p}v_p(s;Z(s))$, the sum running over the plaquettes whose base point lies in $\Box$. The weakened families of Part~IX live on $\mathbb D_{\kappa'}$, $\kappa'=\kappa-C\delta$, and $|Z(s)|_\infty\le gC_T|\xi|_\infty$ there, with $C_T=2^d\sqrt d\,\|C_k^{1/2}\|_{\kappa',\gamma'}$. Indeed, by Lemma~\ref{lem:weak-basic}(a) and Corollary~\ref{cor:weak-covariance}, $|C_k^{1/2}(s)|_{\gamma'}\le2^d\|C_k^{1/2}\|_{\kappa',\gamma'}$. The row sums of the block norms bound the operator on sup norms up to the factor $\sqrt d$ from $\|\xi\|_{B(y')}\le\sqrt d|\xi|_\infty$. By Corollary~\ref{cor:weak-covariance}, $C_T$ is uniform in the reference. Let $|\xi|_\infty\le p$ with $gC_Tp\le\frac12r_I$. Then $V(\Box;s)$ is defined and analytic on the interior of $\mathbb D_{\kappa'}$, continuous up to the boundary. By Corollary~\ref{cor:weak-covariance}, Lemma~\ref{lem:weak-basic}(b) and Proposition~\ref{prop:interaction-pieces}(d), Lemma~\ref{lem:localization} applies to $f(s)=V(\Box;s)$ on $\mathbb D_{\kappa'}$ with $Y=\Box$, and $F=\sup_{\mathbb D_{\kappa'}}|V(\Box;\cdot)|\le C_Pg^{-2}M^d(gC_Tp/r_I)^3$. With $Z(s)=gC_k^{1/2}(s)\xi$, this is of order $g\,|\xi|^3$. The control of $\xi$ belongs to the cluster expansion.

\subsection{The gauge-fixing pieces}\label{subsec:sq-gf}

The exponent of Theorem~\ref{thm:step} contains the remainder $G_3$ of the gauge fixing (Lemma~\ref{lem:exp-gauge}(c)). It is a sum of terms localized in single blocks $B'(y)$. Since $Z(s)$ is not local in $\xi$, it is localized like the interaction. We define its pieces at every reference of the Setting, as Taylor remainders.

Recall $f(H)=\mathsf N-\operatorname{tr}\cos H$, an even entire function with $f(H)=\frac12\operatorname{tr}H^2+O(|H|^4)$ (Section~\ref{sec:fluctuations}). For $y\in\T'$, $x\in B'(y)$ and a complex field $Z$ on $\mathfrak F$, extended by $0$ to the bonds $b_0(c)$, put
\[
H_{y,x}(Z)=\ell\big((e^{\ii Z}W)(y,x)\,W(y,x)^{-1}\big),
\]
at the base $W=U'^{(k)}$ of the reference. It involves only the inner bonds of $B'(y)$, which all lie in $\mathfrak F$.

\begin{definition}[Gauge-fixing pieces]\label{def:gf-pieces}
For complex $Z$ on $\mathfrak F$ with $|Z|_\infty<r_I$ put
\[
\varphi_{y,x}(Z)=f\big(H_{y,x}(Z)\big)-\tfrac12\ip{\mathfrak G_WZ(y,x)}{\mathfrak G_WZ(y,x)},
\]
where $\ip AB=\operatorname{tr}(AB)$ is the bilinear form on $\mathfrak g^c$, and $\mathfrak G_WZ(y,x)=\lambda_W(y,x)[Z]$ (Definition~\ref{def:avg-transport}, in the units of step $k$).
\end{definition}

\begin{proposition}[Gauge-fixing pieces]\label{prop:gf-pieces}
There is $C$, depending only on $d$ and $G$, such that the following holds at every reference of the Setting with $\delta_c<\delta_P$.
\begin{enumerate}[label=(\alph*)]
\item $\varphi_{y,x}$ is defined and analytic for $|Z|_\infty<r_I$, and analytic in $Z$ and in the reference jointly. It does not depend on the weakening parameters $s$.
\item $|\varphi_{y,x}(Z)|\le C(|Z|_\infty/r_I)^3$ for $|Z|_\infty\le\frac12r_I$.
\item $\varphi_{y,x}$ vanishes to third order at $Z=0$: it is $f\circ H_{y,x}$ minus its Taylor polynomial of degree $2$. At real references, $\sum_{x\in B'(y)}\varphi_{y,x}=G_{3,y}$ and $\sum_{y,x}\varphi_{y,x}=G_3$, with $G_{3,y}$, $G_3$ of Lemma~\ref{lem:exp-gauge}(c) at $W_\circ=W$.
\item \emph{(Block locality.)} $\varphi_{y,x}(Z)$ depends on $Z$ only on the inner bonds of $B'(y)$, and on the reference only through the base on these bonds.
\item \emph{(Invariance.)} $\varphi_{y,x}$ is invariant under real gauge transformations $u$, acting by $U'\mapsto U'^u$ and $Z(b)\mapsto R(u(b_-))Z(b)$. It is real at real references for real $Z$.
\end{enumerate}
\end{proposition}

\begin{proof}\hasdetail{prop:gf-pieces}
\emph{Definition and size of $H_{y,x}$.} By the Setting, the real part $W_r=(U'_r)^{(k)}$ of the base lies in $\mathfrak M$. At a real reference $W=W_r$. Since $r_I\le\frac1{32}c_0L^{-1}$, Lemma~\ref{lem:exp-gauge}(c) at $W_\circ=W_r$ shows that $H_{y,x}$ is defined and holomorphic, with $|H_{y,x}(Z)|\le4C_ML|Z|_\infty\le\frac18$; here and below $c_0\le1/(4C_M)$, which is part of the choice ``$c_0$ small'' in Lemma~\ref{lem:exp-gauge}(c). At a complex reference, $W=e^{\ii\beta}W_r$ bondwise, with $|\beta|_\infty\le C_{\rm AB}\delta_c$, and $\beta$ is holomorphic in the complex part of $U'$ (Proposition~\ref{prop:AB}(i), Lemma~\ref{lem:AB-bounds}(a),(b)). By (DC5), $C_{\rm AB}\delta_P\le\frac1{32}c_0L^{-1}$ for $c_{\rm DC}$ small. Write $e^{\ii Z}W=e^{\ii(Z\star\beta)}W_r$ with $Z\star\beta=\ell(e^{\ii Z}e^{\ii\beta})$ bondwise; by Lemma~\ref{lem:log-estimates}(c), $|Z\star\beta|_\infty\le2(|Z|_\infty+|\beta|_\infty)\le\frac18c_0L^{-1}$, since $|Z|_\infty+|\beta|_\infty$ is small. Let $H^r_{y,x}$ be the function of Lemma~\ref{lem:exp-gauge}(c) at $W_\circ=W_r$. Then
\[
H_{y,x}(Z)=\ell\big(e^{\ii H^r_{y,x}(Z\star\beta)}e^{-\ii H^r_{y,x}(\beta)}\big),
\]
with $|H^r_{y,x}(Z\star\beta)|\le\frac18$ and $|H^r_{y,x}(\beta)|\le\frac1{32}$. By Lemma~\ref{lem:log-estimates}(c), $H_{y,x}$ is defined and holomorphic in $(Z,\beta)$, with $|H_{y,x}(Z)|\le C$. By Lemma~\ref{lem:avg-transport}(f), whose hypothesis $L\max|\beta|\le c$ is again part of (DC5), $W$ satisfies the hypothesis of Definition~\ref{def:avg-transport}. So $\mathfrak G_WZ(y,x)=\frac d{dt}\big|_{t=0}H_{y,x}(tZ)$ by Lemma~\ref{lem:avg-transport}(a), and $|\mathfrak G_WZ(y,x)|\le CL\max|Z|$, the maximum over the bonds of $B'(y)$, by Lemma~\ref{lem:avg-transport}(c),(f).

(a) $f$ is entire, and $H_{y,x}$ and $\mathfrak G_W$ are holomorphic in $Z$ and $\beta$, hence in the reference. The definition does not involve $s$.

(c) $H_{y,x}(0)=0$, and the linear part of $H_{y,x}$ is $\mathfrak G_WZ(y,x)$. Since $f(H)=\frac12\operatorname{tr}H^2+O(|H|^4)$, the Taylor polynomial of degree $2$ of $f\circ H_{y,x}$ at $Z=0$ is $\frac12\ip{\mathfrak G_WZ(y,x)}{\mathfrak G_WZ(y,x)}$, at every reference. At real references the sum over $x$ is the definition of $G_{3,y}$ in the proof of Lemma~\ref{lem:exp-gauge}(c).

(b) Fix $Z$ with $0<|Z|_\infty\le\frac12r_I$, and put $\phi(t)=\varphi_{y,x}(tZ)$ for $|t|<r_I/|Z|_\infty$. There $|H_{y,x}(tZ)|\le C$ and $|\mathfrak G_W(tZ)(y,x)|\le CLr_I\le C$, so $|\phi(t)|\le C$, since $f$ is entire. By (c), $\phi(t)/t^3$ is holomorphic on this disc, and the maximum principle gives $|\varphi_{y,x}(Z)|=|\phi(1)|\le C(|Z|_\infty/r_I)^3$.

(d) The averaged contour variable $(e^{\ii Z}W)(y,x)$ depends only on the bonds of the box spanned by $y$ and $x$, which are inner bonds of $B'(y)$ (Section~\ref{sec:fluctuations}); so does $\lambda_W(y,x)$ (Lemma~\ref{lem:avg-transport}(a)).

(e) For a real gauge transformation $u$ the base becomes $W^u$ (Setting), and $e^{\ii R(u)Z}W^u=(e^{\ii Z}W)^u$ bondwise. By \eqref{eq:avg-contour-gauge}, which holds for real $u$ whenever the averaged contour variables are defined, $H_{y,x}$ becomes $\ell\big(u(y)e^{\ii H_{y,x}}u(y)^{-1}\big)=u(y)H_{y,x}u(y)^{-1}$ (Lemma~\ref{lem:log-estimates}(d)). By Lemma~\ref{lem:avg-transport}(d), $\mathfrak G_WZ(y,x)$ becomes $u(y)\,\mathfrak G_WZ(y,x)\,u(y)^{-1}$. Both terms of $\varphi_{y,x}$ are invariant under this conjugation. At a real reference and for real $Z$, $H_{y,x}$ and $\mathfrak G_WZ(y,x)$ lie in $\mathfrak g$, so $\varphi_{y,x}$ is real.
\end{proof}

The pieces are grouped by boxes like the interaction. Every $y\in\T'$ lies in exactly one box. Put
\[
G_3(\Box;s)=g^{-2}\sum_{y\in\T'\cap\Box}\ \sum_{x\in B'(y)}\varphi_{y,x}\big(Z(s)\big).
\]
At real references, $\sum_\Box G_3(\Box;1)=g^{-2}G_3(Z)$ for $Z=gC_k^{1/2}\xi$, by Proposition~\ref{prop:gf-pieces}(c) and $C_k^{1/2}(1)=C_k^{1/2}$ (Corollary~\ref{cor:weak-covariance}). Let $|\xi|_\infty\le p$ with $gC_Tp\le\frac12r_I$. Then $G_3(\Box;s)$ is defined and analytic on the interior of $\mathbb D_{\kappa'}$, continuous up to the boundary. A box contains at most $(2M)^d$ pairs $(y,x)$, so by Proposition~\ref{prop:gf-pieces}(b)
\[
\sup_{\mathbb D_{\kappa'}}|G_3(\Box;\cdot)|\le C(2M)^dg^{-2}\big(gC_Tp/r_I\big)^3\le C_Pg\,p^3 .
\]
The terms of $G_3(\Box;s)$ are local in $Z$, with range $L$, but $Z(s)$ is not local in $\xi$. As for $\mathfrak J(Z(s))$ below, Lemma~\ref{lem:localization} applies to $f(s)=G_3(\Box;s)$ on $\mathbb D_{\kappa'}$ with $Y=\Box$, by Corollary~\ref{cor:weak-covariance}, Lemma~\ref{lem:weak-basic}(b) and Proposition~\ref{prop:gf-pieces}(d); the blocks $B'(y)$, $y\in\Box$, lie within distance $L$ of $\Box$, which the margin $r'=r+2L+2$ covers. The pieces are analytic in the reference and invariant under real gauge transformations, by Proposition~\ref{prop:gf-pieces}(a),(e) and Lemma~\ref{lem:localization}(c).

\subsection{The Jacobians}\label{subsec:sq-jacobians}

\begin{proposition}[Jacobians]\label{prop:jacobians}
\begin{enumerate}[label=(\alph*)]
\item $\mathfrak J(Z)=\sum_{b\in\mathfrak F}\log(\sigma(Z(b))/\sigma(0))+\sum_c\log(\mathsf J_c(Z)/\mathsf J_c(0))$ is a sum of terms, each analytic for $|Z|_\infty\le r_L$, depending on $Z$ and on the base only within range $2L$ of one bond, vanishing at $Z=0$ and bounded as in Theorem~\ref{thm:step}(e). The terms are invariant under real gauge transformations. Since Proposition~\ref{prop:exact-step} is written with $\det$, which is positive at real references, they are analytic in $Z$ and in the reference jointly, on the complex references with $\delta_c<\delta_P$, and real at real references and real $Z$.
\item By Proposition~\ref{prop:exact-step}, $\mathfrak j_c=\det(1+\mathfrak r_c(0))^{-1}$, with $\mathfrak r_c$ from Lemma~\ref{lem:eliminate}(b). It is an analytic function of the base near $c$, invariant under real gauge transformations, and $\mathfrak j_c=1$ at the trivial base.
\item Let $G$ be semisimple. As a function of $U'$ (in the units of step $k$) through $W=U'^{(k)}$, $\log\mathfrak j_c$ is a normalized graded activity at the blocks $B'(c_-)\cup B'(c_+)$ on $\mathcal U(\cdot;c_a,c_a,c_a)$ for $c_a\le c_JL^{-d-2}$ ($=c_JL^{-6}$ in $d=4$), with sup at most $C$. Hence $|\log\mathfrak j_c|\le C$ on $\mathcal U(\cdot;\theta_r,\delta_1,\delta_0)$ for $\theta_r,\delta_1,\delta_0\le c_a$. Moreover, with $\mathfrak t_0$ as in Proposition~\ref{prop:one-loop-polymer}(d), $|\log\mathfrak j_c|\le C_P\mathfrak t_0^2$ whenever $C_P(\mathfrak t_0+\delta_0)\le1$. (For abelian $G$, $\mathfrak r_c\equiv0$ and $\mathfrak j_c\equiv1$.)
\item \emph{(Lattice covariance.)} Let the base $W$ be real, as in Section~\ref{sec:fluctuations}, and let $r\in\mathfrak S_{k+1}$. Then $rW$ is the base of $rV'$ for the background field $rU'$ (Proposition~\ref{prop:symmetry}(b)), and $\mathfrak j_{c^r}(rW)=\mathfrak j_c(W)$ for every positively oriented bond $c$ of $\T'$, where $c^r$ is the positively oriented bond with $c^r=\pm rc$. Hence $\sum_c\log\mathfrak j_c(rW)=\sum_c\log\mathfrak j_c(W)$. The same holds at $k=0$, with $\mathfrak S_1$ and the base $W=U'$.
\end{enumerate}
\end{proposition}
Part~(d) is used only in Part~XI, in the proof of Lemma~\ref{lem:ol-symmetry}.

\begin{proof}\hasdetail{prop:jacobians}
(a) is Proposition~\ref{prop:exact-step} and Theorem~\ref{thm:step}(e), with Lemma~\ref{lem:eliminate}(a),(c) for the locality. For invariance, a real gauge transformation $u$ acts on the relative variables by $Z(b)\mapsto R(u(b_-))Z(b)$ and on $\ell_c$ by $R(u(c_-))$. The function $\sigma$ is $\operatorname{Ad}$-invariant, and $\det\partial\ell_c/\partial Z(b_0(c))$ changes by $\det R(u(c_-))\det R(u(b_0(c)_-))^{-1}=1$. The terms are analytic functions of $Z$ and of the base, through $\ell_c$ and $\zeta_c$, and the base and the elimination are analytic in the reference (Setting). At real references with $L^d\mathsf a\le c$ and $|Z|_\infty\le r_L$, $\partial\ell_c/\partial Z(b_0(c))$ is real, and its determinant is positive (Proposition~\ref{prop:exact-step}). So the analytic terms are those of Theorem~\ref{thm:step}(e) there.

(b) At $Z=0$ we have $\zeta_c(0)=0$, since $\ell_c$ vanishes at $Z=0$. By Lemma~\ref{lem:eliminate}(b) and $\det R=1$ on $\mathfrak g$, $\det\partial\ell_c/\partial Z(b_0(c))=L^{(1-d)\dim\mathfrak g}\det(1+\mathfrak r_c(0))$. This gives the formula. At the trivial base, $l^{\rm pl}=q^{\rm pl}$ of \eqref{eq:q-plain}, whose coefficient of $Z(b_0(c))$ is $L^{1-d}$ exactly. So $\partial\ell_c/\partial Z(b_0(c))=L^{1-d}$ and $\mathfrak r_c(0)=0$; in $d=4$ the coefficient is $L^{-3}$. Invariance is (a) at $Z=0$. For complex bases, $\det R(g)=1$ for $g\in G^c$, since $G$ is connected (Section~\ref{sec:background}). So the formula is the analytic continuation.

(c) On $\mathcal U(P;c_a,c_a,c_a)$, the base $W=U'^{(k)}$ is a complex perturbation of size $Cc_a$ of a real base of curvature at most $Cc_a$, in the units of step $k$ (Proposition~\ref{prop:avg-regular}, Proposition~\ref{prop:AB}(i)). In Lemma~\ref{lem:eliminate}, $\mathsf a$ is a curvature in the units of step $k+1$ (Setting), which is $L^2$ times the curvature in the units of step $k$. So $L^d\mathsf a\le CL^{d+2}c_a$ for the real part. With the plain constraint, $l^{\rm pl}=Lp'$ and $p'=l'-D_{W'}\gamma'$ (Lemma~\ref{lem:eliminate}(a), Lemma~\ref{lem:one-step-expansion}(f)). The contour average $\gamma'$ involves only inner bonds, so it does not involve $Z(b_0(c))$, and the coefficient of $Z(b_0(c))$ comes from $l'$ alone. The complex part of the base contributes to $\mathfrak r_c(0)$ through the remainder $r_j$ of Proposition~\ref{prop:AB}(ii) for the step from scale $k$ to scale $k+1$. Proposition~\ref{prop:AB} is stated for $j\le k$; we apply it with $k+1$ in place of $k$, in the units of step $k+1$, where the curvature becomes $L^2Cc_a$ and the size of the complex part $LCc_a$. Its hypotheses hold for $c_J$ small: $L^2Cc_a\le c'_3$ and $LCc_a<\delta_{AB}$. The remainder is relatively $CL^{d+1}c_a$. Hence $|\mathfrak r_c(0)|\le CL^{d+2}c_a\le\frac12$ for $c_J$ small. Then $\log\mathfrak j_c=-\log\det(1+\mathfrak r_c(0))$ is analytic and bounded by $C$. It is normalized by (b), invariant by (b), and depends on $U'$ only on the bonds of $B'(c_-)\cup B'(c_+)$. The first bound on the smaller domains follows since the domains are monotone in their constants. For the second, apply Lemma~\ref{lem:quadratic-bound} with $\theta_r=\delta_1=\delta_0=c_a$ and a cube $\Lam$ of bounded side around the two blocks, as in the proof of Proposition~\ref{prop:one-loop-polymer}(c). Then $\mathfrak s\le C_P\mathfrak t$, and $C_P(\mathfrak t_0+\delta_0)\le1$ gives the hypotheses of the lemma and $\mathfrak s\le\frac12$.

(d) By Proposition~\ref{prop:symmetry}(b), $rU'$ is an element of the minimal orbit of $rV'$, and its base is $(rU')^{(k)}=rW$, by Lemma~\ref{lem:one-step-basic}. Write $\ell_c(Z;W)$ for $\ell_c$ at the base $W$, and $(r_WZ)(b)=Z(r^{-1}b)$, with $Z(-b)=-R(W(b))^{-1}Z(b)$, as in Proposition~\ref{prop:hessian}(e). By the convention on reversed bonds, $r(e^{\ii Z}W)=e^{\ii r_WZ}\,rW$ bondwise. By Lemma~\ref{lem:one-step-basic}(d) and Lemma~\ref{lem:reversal}, $\mathfrak a'(r(e^{\ii Z}W))=r\,\mathfrak a'(e^{\ii Z}W)$; in particular $\mathfrak a'(rW)=rW'$. So the relative average at the base $rW$ is $\mathfrak a'(e^{\ii r_WZ}rW)\,(rW')^{-1}=r\overline X$, with $\overline X=\mathfrak a'(e^{\ii Z}W)W'^{-1}$ as in Section~\ref{sec:fluctuations} and $(r\overline X)(c')=\overline X(r^{-1}c')$. Let $c$ be a positively oriented bond of $\T'$, and $c^r=\epsilon\,rc$ with $\epsilon=\pm1$. If $\epsilon=1$, $(r\overline X)(c^r)=\overline X(c)$ and $\ell_{c^r}(r_WZ;rW)=\ell_c(Z;W)$. If $\epsilon=-1$, Lemma~\ref{lem:reversal} gives $(r\overline X)(c^r)=\mathfrak a'(e^{\ii Z}W)(-c)\,W'(-c)^{-1}=R(W'(c))^{-1}\big(\overline X(c)^{-1}\big)$, so $\ell_{c^r}(r_WZ;rW)=-R(W'(c))^{-1}\ell_c(Z;W)$, by Lemma~\ref{lem:log-estimates}(d). So $\ell_{c^r}(r_WZ;rW)=O_c\,\ell_c(Z;W)$ with an orthogonal map $O_c$ of $\mathfrak g$, since $W'$ is real. The map $r$ sends $b_0(c)$ onto $\pm b_0(c^r)$ (Section~\ref{sec:fluctuations}), so $(r_WZ)(b_0(c^r))=O'_cZ(b_0(c))$ with $O'_c=1$ or $O'_c=-R(W(b_0(c)))^{-1}$, also orthogonal, and $Z(b_0(c))$ enters $r_WZ$ only there. Differentiating at $Z=0$ with respect to $Z(b_0(c))$,
\[
\frac{\partial\ell_{c^r}}{\partial Z(b_0(c^r))}(0;rW)\,O'_c=O_c\,\frac{\partial\ell_c}{\partial Z(b_0(c))}(0;W).
\]
Orthogonal maps have determinant $\pm1$, so the two determinants agree up to sign. Both are positive at the real base (Proposition~\ref{prop:exact-step}), so they are equal. Since $\zeta_c(0)=0$, Proposition~\ref{prop:exact-step} gives $\mathfrak j_c=L^{-(d-1)\dim\mathfrak g}\mathsf J_c(0)=L^{-(d-1)\dim\mathfrak g}\,\big(\det\partial\ell_c/\partial Z(b_0(c))\big)^{-1}$ at $Z=0$, and the same at $(c^r,rW)$. So $\mathfrak j_{c^r}(rW)=\mathfrak j_c(W)$. As $c$ runs over the positively oriented bonds of $\T'$, so does $c^r$, which gives the sum. At $k=0$ the same argument applies: Lemma~\ref{lem:one-step-basic}(d), Lemma~\ref{lem:reversal}, Lemma~\ref{lem:eliminate} and Proposition~\ref{prop:exact-step} are one-step statements that hold at $k=0$ (proof of Proposition~\ref{prop:first-step}(a)), and $rU'$ is an element of the minimal orbit of $rV'$ by Theorem~\ref{thm:balaban}(c) for the one-step problem.
\end{proof}

Part~(d) uses only the intrinsic objects of the step: the averaging, the bonds $b_0(c)$ and the elimination at a real base. It is used for the lattice symmetry of the one-loop term of the step at the background fields (Lemma~\ref{lem:ol-symmetry}(b)).

The terms of $\mathfrak J(Z(s))$ are local in $Z$, but $Z(s)$ is not local in $\xi$. So $\mathfrak J(Z(s))$, grouped by boxes, goes through Lemma~\ref{lem:localization} like $V(\Box;s)$.

\subsection{The composed activities}\label{subsec:sq-composed}

The norm $|\cdot|_*$ of Theorem~\ref{thm:weak-bg} and the quantities of Definition~\ref{def:graded} are compared as follows.

\begin{lemma}[Norm conversion]\label{lem:norm-conversion}
There is $C_{\rm conv}$, depending only on $d$, $\alpha$ and $G$ (through Theorem~\ref{thm:ID}(b)), with the following property. Let $\alpha\le\frac12$, so that the exponent $\beta=\min(\alpha,\frac12)$ of $|\cdot|_*$ equals $\alpha$. Let $U'_r$ be real and $(M,\eps)$-regular with $\eps\le\eps_0$, and let $A$ be a $\mathfrak g^c$-valued field on the bonds of $\Tf$, with $|A|_*$ taken in the real gauges of $U'_r$ (Section~\ref{sec:weak-background}). Then for every set $Y$,
\[
|A|_Y+|\nabla A|_Y+[\nabla A]_{\alpha,Y}\le C_{\rm conv}|A|_*,
\]
with $|A|_Y$ the sup norm on the bonds of $Y$, and $\nabla$ and the H\"older quotient of Definition~\ref{def:graded} taken with respect to $U'_r$.
\end{lemma}

\begin{proof}\hasdetail{lem:norm-conversion}
Every bond lies in a block $B(y)$, and $\mathfrak n^3_y(A)$ bounds $|A|_{B(y)}$ and $|\nabla A|_{B(y)}$. If the gradient in $\mathfrak n^3_y$ is read in the gauge of $y$ as a flat derivative, \eqref{eq:cov-vs-flat} with $\sigma=C\eps$ converts it, with an error $C\eps|A|$. Consider a quotient in $[\nabla A]_{\alpha,Y}$, with $x\ne x'$ in a cube $\Lam_3(z')\subset Y$. If $|x-x'|\ge1$, it is at most $|\nabla A(x)|+|\nabla A(x')|\le2|\nabla A|_Y$, since $R(U'_r(\Gamma_{x,x'}))$ is orthogonal. If $|x-x'|<1$, then $x,x'$ lie in a common cube $\Lam_1(c_y)$ (Section~\ref{sec:weak-background}, setting). In the real gauge of $y$, $U'_r=e^{\ii\eta A_r}$ with $\|A_r\|_{1,\alpha}\le C\eps$ on $\Lam_4(c_y)$ (Theorem~\ref{thm:ID}(b)). There \eqref{eq:cov-vs-flat} and $\|R(U'_r(\Gamma_{x,x'}))-1\|\le C|x-x'|\eps$ bound the quotient by $[\partial\tilde A_y]_{\alpha,\Lam_1(c_y)}+C\eps(|A|+|\partial\tilde A_y|)\le C\mathfrak n^3_y(A)$. Taking the supremum over $y$ gives the bound.
\end{proof}

In $C_{\rm conv}$, the pairs at distance at least $1$ cost a factor $2$.

\begin{proposition}[Composed activities]\label{prop:composed}
Let $X$ be a connected domain of step $k$ (an old polymer), let $r_E\ge0$ with $r_E+3\le r$, and put $P=X^{(r_E)}$. Let $\mathcal E(X,\cdot)$ be a graded activity at $P$ on $\mathcal U(P;\theta_r,\delta_1,\delta_0)$ (Definition~\ref{def:graded}), with sup $\|\mathcal E(X)\|$; so it depends on the configuration on $P$, and its domain is controlled on $P^{(3)}=X^{(r_E+3)}$. Let $U'=e^{\ii\eta A'}U'_r$ be a reference as in the Setting, and suppose that on $P^{(3)}$ this is a representation as in Definition~\ref{def:graded}, with budgets $\theta_r,\delta'_1,\delta'_0$, where $\delta'_1<\delta_1$ and $\delta'_0<\delta_0$. Thus the real part of the representation is the $U'_r$ in whose gauges $|\cdot|_*$ is measured. Let $0<\zeta\le\frac12R_P$, and put
\[
\tau=C_1C_{\rm conv}C_P\,\zeta\,\Big(\min(\delta_1-\delta'_1,\ \delta_0-\delta'_0)\Big)^{-1}.
\]
Then, for $\tau<\frac12$, $s\in\mathbb D_\kappa$ and $|\hat Z(Z)|_\infty\le\zeta$:
\begin{enumerate}[label=(\alph*)]
\item $e^{\ii\eta A_*(s;\hat Z(Z))}U'\in\mathcal U(P;\theta_r,\delta_1,\delta_0)$, and
\[
f(X;s,Z)=\mathcal E\big(X,e^{\ii\eta A_*(s;\hat Z(Z))}U'\big)-\mathcal E(X,U')
\]
is analytic in $(s,Z)$ and in the reference jointly, continuous up to the boundary in $s$, with $|f(X;s,Z)|\le4\tau\|\mathcal E(X)\|$;
\item if $s_\Box=0$ outside a union of boxes $Z'$ that contains the boxes meeting $X$, then $f(X;s,Z)$ depends on $s$ only through the boxes of the $*$-component $Z'_0$ of $Z'$ containing them, and on $Z$ and $U'$ only on $(Z'_0)^{(r')}$;
\item $f(X;s,Z)$ is invariant under real gauge transformations.
\end{enumerate}
\end{proposition}

\begin{proof}\hasdetail{prop:composed}
(a) For $|t|<\frac1{2\tau}$ consider $e^{\ii\eta tA_*}U'=e^{\ii\eta(tA_*\oplus A')}U'_r$ on $P^{(3)}$, with $A_*=A_*(s;\hat Z(Z))$, combined bondwise by the Baker--Campbell--Hausdorff formula. The real part $U'_r$ is unchanged, so no current of $A_*$ is needed. By Theorem~\ref{thm:weak-bg}(a), $|A_*|_*\le C_P\zeta$. By Lemma~\ref{lem:norm-conversion}, the sup norm, the covariant gradient and the H\"older quotient of $A_*$ on $P^{(3)}$, with respect to $U'_r$, are at most $C_{\rm conv}C_P\zeta$. By the Baker--Campbell--Hausdorff formula and \eqref{eq:lattice-trick}, the corresponding quantities of $tA_*\oplus A'$ exceed those of $A'$ by at most $C_1C_{\rm conv}C_P|t|\zeta<\frac12\min(\delta_1-\delta'_1,\delta_0-\delta'_0)$, for a suitable $C_1$. So $tA_*\oplus A'$ satisfies the bounds of $\mathcal U(P;\theta_r,\delta_1,\delta_0)$ on $P^{(3)}$, strictly. Hence $t\mapsto\mathcal E(X,e^{\ii\eta tA_*}U')-\mathcal E(X,U')$ is analytic for $|t|<\frac1{2\tau}$, bounded by $2\|\mathcal E(X)\|$, and vanishes at $t=0$. Since $\tau<\frac12$, the point $t=1$ lies in the disc, and the Schwarz lemma gives $|f|\le2\|\mathcal E(X)\|\cdot2\tau$ there.

Analyticity. $A_*(s;B)$ is analytic in $(s,B)$ and in $A'$ jointly (Theorem~\ref{thm:weak-bg}(b), Proposition~\ref{prop:complex-ext}(d)). $\hat Z$ is analytic in $Z$ and in the base, and the base is analytic in $A'$ (Proposition~\ref{prop:AB}(vi), Setting). $\mathcal E(X,\cdot)$ is holomorphic on an open set containing $\mathcal U(P)|_P$, and the bounds above hold strictly. So $f$ is analytic in $(s,Z,A')$ jointly, as a composition. By (b), it depends on $A'$ only on $(Z'_0)^{(r')}$, so it is analytic in the reference near that set. Continuity up to the boundary in $s$ follows from Theorem~\ref{thm:weak-bg}(b).

(b) $\mathcal E(X,\cdot)$ depends on the configuration on $P=X^{(r_E)}\subset(Z'_0)^{(r)}$, and $A_*$ on $(Z'_0)^{(r)}$ is determined by the data on $(Z'_0)^{(r')}$ and by $s$ in $Z'_0$ (Theorem~\ref{thm:weak-bg}(d), and Lemma~\ref{lem:eliminate} for $\hat Z$). (c) follows from the covariance of $A_*$ and $\hat Z$ and the invariance of $\mathcal E$.
\end{proof}

In the step, $\mathcal E=\mathcal E_k$ on the graded domain $\mathcal U^\flat(X)=\mathcal U(X;\theta_r,\delta_1^\flat,\delta_0)$ of step $k$, $\delta_1^\flat=\frac14\delta_1$, on which the activities carried by the flow are bounded (Section~\ref{sec:generations}); the proposition is applied with $\delta_1^\flat$ in place of $\delta_1$, with $r_E=0$ and $P=X$, and $U'$ is the extension (Setting) of a point of the domain of step $k+1$, in the units of step $k$. On $X^{(3)}$ the extension is not felt, since the pieces used have $Z^{(r')}$ inside the polymer of step $k+1$ (Setting), and $X^{(3)}\subset Z^{(r')}$. By Remark~\ref{rem:domain-constants}(i), the representation on $X^{(3)}$ has the same real part and budgets $\delta'_1=L^{-2}\delta_1$, $\delta'_0=L^{-1}\delta_0$. So $\min(\delta_1^\flat-\delta'_1,\delta_0-\delta'_0)\ge(\frac14-\frac19)\delta_1\ge\frac18\delta_1$, since $L\ge3$ and $\delta_1\le\delta_0$ (Remark~\ref{rem:symmetrization-flow}(c)). On the support of the small-field cutoff, $\zeta$ is of order $L^dg_kp_{0,k}$ ($\zeta=w$, Section~\ref{sec:cluster}). So $\tau\le C_PL^dg_kp_{0,k}/\delta_1$, which is small when $g_kp_{0,k}$ is, the domain constants being fixed. For $|\xi|_\infty\le p$ with $CL^dgC_Tp\le\zeta$, the function $s\mapsto f(X;s,Z(s))$ is analytic on the interior of $\mathbb D_{\kappa'}$ and bounded by $4\tau\|\mathcal E(X)\|$. Lemma~\ref{lem:localization} applies to it on $\mathbb D_{\kappa'}$, without conversion of margins (Remark~\ref{rem:localization-margins}). Here $Y$ is the set of boxes of $X$ when $X$ is a union of boxes, as for the polymer grid; in general it is the union of the boxes meeting $X$.

The next remark uses Theorem~\ref{thm:mayer}(a), the enlarged closure $\hat Y$ of Remark~\ref{rem:reblock-margins} and, in its last paragraph, the grouped old activities of Definition~\ref{def:old-activities}. It may be read after Section~\ref{subsec:flow-reblock}.

\begin{remark}[Extended references]\label{rem:extension}
A point of the domain of step $k+1$ at a polymer $Y'$ fixes $U'$ only near $Y'$. It is extended by the cutoff of Remark~\ref{rem:bg-to-graded}. In the cutoff band, $|\nabla(\chi A)|\sim|A|$, which need not be within the gradient budget $\delta_1$ of the domain of step $k$. So composed configurations of old polymers meeting the band need not lie in the domain of step $k$ at $X$. Therefore, before the Mayer expansion, all old activities and all pieces whose polymer is not contained in $Y'$ (in the units of step $k$) are set to zero. By Theorem~\ref{thm:mayer}(a), the output $E^\#(Y)$ with $\hat Y=Y'$ is unchanged, since it depends only on the pieces of polymers inside $Y$. So the new activity at $Y'$ is defined at every point of the domain of step $k+1$ at $Y'$ and does not depend on the extension. Pieces with $Z\subset Y'$ near $\partial Y'$ may depend on the extension, but they do not enter $E^\#(Y)$ with $\hat Y=Y'$. At a real critical $U'$, take $U_r=U'$ and $A=0$, so that no extension is needed. The restricted and the full systems then share every $E(Z)$ with $Z\subset Y$. By Theorem~\ref{thm:mayer}(a), applied also to $E^\#_{\rm p}$ and to the $\lambda$-family of Corollary~\ref{cor:two-classes}, their outputs agree. Since $\sum_{Y'}\sum_{\hat Y=Y'}=\sum_Y$, the identity of Theorem~\ref{thm:sf-step} is not affected.

\emph{The collar inclusions.} Let $Y$ contribute to the output at $Y'$, i.e.\ $\hat Y=Y'$, with $\hat Y$ the enlarged closure of Remark~\ref{rem:reblock-margins}. Then $Y^{(r')}\subset Y'$, and the following holds.
\begin{itemize}
\item By Theorem~\ref{thm:mayer}(a), the output at $Y$ uses only pieces $E(Z)$ with $Z\subset Y$.
\item Each such piece depends on the reference only on $Z^{(r')}\subset Y^{(r')}\subset Y'$.
\item A mixed piece uses old activities $\mathcal E_k(X)$ with $X\subset Z$. Their domain budget is needed on $X^{(3)}\subset Z^{(r')}\subset Y'$, since $r'=r+2L+2\ge3$.
\end{itemize}
So on every set that matters the extension is not felt. The real part there is that of $U'$, and the slack of Remark~\ref{rem:domain-constants}(i) gives the hypotheses of Proposition~\ref{prop:composed}. The same inclusions hold for the pure system and for the $\lambda$-family of Corollary~\ref{cor:two-classes}.

To avoid an implicit definition, the restricted system for an output polymer $Y'$ is defined on the input pieces with $Z\subset Y'$ only. Equivalently, it may be defined on a cutoff extension that equals $U'$ on the $3$-neighbourhood of $Y'$, which is the margin of the next domain. Pieces near $\partial Y'$ whose collars cross it may depend on the extension, but they do not enter the outputs with $\hat Y=Y'$. At real critical references the restricted and the unrestricted systems share all inputs that enter, so their outputs coincide.

\emph{Grouped old activities.} In the flow the old activities of step $k$ are the grouped activities $\mathcal E_k(Y)$ of Definition~\ref{def:old-activities}, indexed by the polymers $Y$ of scale $k$, the polymers of the paving. The activity $\Psi^{(j)}_\Box$ of a grid cube $\Box$, which contains the counterterm $-\beta_jS_\Box$, is carried by the polymer $\Box^+$. Each $\mathcal E_k(Y)$ is a graded activity at $Y$ itself, on $\mathcal U^\flat(Y)$, with $r_E=0$ and $P=Y$ (Proposition~\ref{prop:old-activities}). So everything above applies with the old polymer $X=Y$. The restricted system for $Y'$ contains $\mathcal E_k(Y)$ if and only if $Y\subset Y'$, and $\mathcal E_k(Y)$ is set to zero or kept as a whole. In particular the counterterm $-\beta_jS_\Box$ is never separated from the small polymers of $\Box$. A mixed piece that uses $\mathcal E_k(\Box^+)$ has $\Box^+\subset Z$, and the budget of $\mathcal E_k(\Box^+)$ is needed on $(\Box^+)^{(3)}\subset Z^{(r')}\subset Y'$, as in the collar inclusions. The single activities $E^{(j)}(X,\cdot)$ inside $\mathcal E_k(Y)$ are functions of the configuration on $X\subset Y$, and their domains are controlled on $X^{(3)}\subset Y^{(3)}$ (Remark~\ref{rem:domain-constants}(i)), so they need no extension of their own. The new activity at $Y'$ is the raw new generation of Definition~\ref{def:raw-new}.
\end{remark}

\subsection{The Hessian away from critical backgrounds}\label{subsec:sq-hessian}

Section~\ref{sec:localization} already defines $\Delta_k=\Delta_k^{(\Tf)}$ and $C_k=\Delta_k^{-1}$ at regular backgrounds that need not be critical, with the elimination at the base $U^{(k)}$ (Definition~\ref{def:dirichlet-hessian} and the definition following it). We use this definition at every reference. At complex references, $J_U$ is not bounded uniformly in $\eta$. So $\omega_0=\bH_0^TJ_U$ is continued in the form
\[
\omega_0=(D\bH_0)^TF^s_{U'},
\]
which agrees with $\bH_0^TJ_U$ at real references, since $J_U=D^TF^s_U$ (Lemma~\ref{lem:current}). Here $\bH_0$ is $\bH$ with potential $0$, and $F^s_{U'}$ is continued holomorphically. This is the form used at level~(L4) of Section~\ref{sec:nested-walks}.

\begin{proposition}[Hessian at a general reference]\label{prop:hessian-general}
\begin{enumerate}[label=(\alph*)]
\item At real critical references, $\Delta_k$ agrees with Lemma~\ref{lem:step-quadratic} and Corollary~\ref{cor:covariance}, including the gauge-fixing term $\ip{\mathfrak G_W\cdot}{\mathfrak G_W\cdot}$.
\item $|\omega_0|_\infty\le C(\eps+\delta_c)$. The operator $W_2$ is real symmetric at real references, with $\|W_2\|\le C(\eps+\delta_c)$ and $L^2$ decay. Only the $L^2$ theory of Part~II is used for it, not the pointwise class $\cW(\eps,\gamma)$.
\item $\Delta_k$ has decay $(C_L,\gamma)$ for some $\gamma>0$.
\item For $\eps\le\eps_{\max}$ and $\delta_c<\delta_P$, so that $\eps+\delta_c$ is small depending on $L$ (Remark~\ref{rem:domain-constants}), $\Rea\Delta_k(Z,\bar Z)\ge\frac18c_L\|Z\|^2$.
\item $\Delta_k$ is analytic in the reference and covariant under real gauge transformations. At $s=1$ it is represented by the weakened family $\Delta^{\rm c}_{\Tf}(s)$ of level~(L4), plus the gauge-fixing part $\mathfrak G_W^T\mathfrak G_W$, which is not weakened.
\end{enumerate}
\end{proposition}

\begin{proof}\hasdetail{prop:hessian-general}
(a) is Lemma~\ref{lem:hessian-gamma} and Proposition~\ref{prop:dirichlet-hessian}(a), together with $K^{\Tf}=\mathfrak G_W^T\mathfrak G_W$ (Section~\ref{sec:localization}).

(b) At real references this is Proposition~\ref{prop:dirichlet-hessian}(b) and Lemma~\ref{lem:hessian-gamma}. At complex references use the block $L^2$ decay of $D\bH_0$ and $|F^s_{U'}|\le C(\eps+\delta_c)$ (Proposition~\ref{prop:complex-ext}(c)). The bound $\|W_2\|\le C(\eps+\delta_c)$ and the $L^2$ decay at complex references come from Proposition~\ref{prop:dirichlet-hessian}(c). There $d^2\mathcal N(0)$ is bounded through its plaquette form, with the bounds $|R(H_p)-1|\le C\eta^2(\eps+\delta_c)$ and $|F^s|\le C(\eps+\delta_c)$ of step~(v) of the proof of Proposition~\ref{prop:complex-ext}, which use the derivative part of $\delta_c$. The $L^2$ operators at complex references come from Proposition~\ref{prop:reg-complex}, through step~(i) of the proof of Proposition~\ref{prop:complex-ext}.

(c) is Proposition~\ref{prop:dirichlet-hessian}(b) for $X=\Tf$, together with the locality of $\hat Z^1$, $\hat Z^2$ and the block-locality of $\mathfrak G_W$. At complex references the coefficient operators (the averages and the covariant derivatives) change by local operators of size $C_L\delta_c$ (Proposition~\ref{prop:reg-complex}, Proposition~\ref{prop:AB}(v)), and by Lemma~\ref{lem:perturbed-inverse} the inverses in $\Delta_k$ change by $C_L\delta_c$ in the decay norm, and $\mathfrak G_W$ by a block-local operator of the same size (Lemma~\ref{lem:avg-transport}(f) with (DC5)), which preserves the decay.

(d) This is Lemma~\ref{lem:coercive-general}, which does not use criticality, applied at the real part $U'_r$; it needs $|J_{U'_r}|\le\eps$ (Setting); its conditions $\eps\le\eps_{\rm cc}$ and those of the proof of Proposition~\ref{prop:dirichlet-hessian}(c) hold since $\eps\le\eps_{\max}$ (Remark~\ref{rem:domain-constants}). The lemma covers complex perturbations of size at most $c''_L$, and a complex reference with $\delta_c<\delta_P$ is one, since $\delta_P\le c''_L$ (Proposition~\ref{prop:complex-ext}, smallness). As required by Proposition~\ref{prop:local-fluct}(b), $\omega_0$ is read there in the form $(D\bH_0)^TF^s_{U'}$.

(e) Analyticity and covariance are Proposition~\ref{prop:dirichlet-hessian}(c), and for the gauge-fixing part Lemma~\ref{lem:avg-transport}(a),(d),(f). Away from criticality, covariance follows directly from the covariance of $\boldsymbol\Gamma_\pi[W_2]$, $\bH_0$, $F^s$, $\hat Z$ and $\mathfrak G_W$. The representation is the construction of level~(L4) with $W_2$, and the gauge-fixing part as in level~(L5).
\end{proof}

\begin{remark}[Source-free step]\label{rem:source-free}
For a general reference, the step quantities are those defined in this section. The exponent is $-g^{-2}\sum_pv_p(s;Z(s))-g^{-2}\sum_{y,x}\varphi_{y,x}(Z(s))$, with no terms of order $0$, $1$ or $2$ in $Z$. The Gaussian integral is taken through the white-noise representation $Z=gC_k^{1/2}\xi$, with $C_k=\Delta_k^{-1}$ as in Section~\ref{subsec:sq-hessian}; at complex references this representation, not a Gaussian measure, defines it. The determinant densities are those of Proposition~\ref{prop:one-loop-polymer}, the Jacobians those of Proposition~\ref{prop:jacobians}, and the composed activities those of Proposition~\ref{prop:composed}. The gauge-fixing pieces are those of Proposition~\ref{prop:gf-pieces}. At real critical references they reproduce Theorem~\ref{thm:step}, by Propositions~\ref{prop:interaction-pieces}(c), \ref{prop:gf-pieces}(c) and~\ref{prop:hessian-general}(a), and Theorem~\ref{thm:weak-bg}(c) for the composed activities: there $\sum_pT_2[\cdots]+\frac12\ip{\mathfrak G_WZ}{\mathfrak G_WZ}=S(U')+\frac12\Delta_k(Z,Z)$. At a non-critical reference this identity fails. This is harmless, since the exponent is a remainder by definition.
\end{remark}

\begin{remark}[What the cluster expansion receives]\label{rem:step-quantities-output}
\begin{enumerate}[label=(\alph*)]
\item \emph{Pieces.} Lemma~\ref{lem:localization} applied to $V(\Box;s)$, to $G_3(\Box;s)$, to the boxes of $\mathfrak J(Z(s))$, and to $f(X;s,Z(s))$ gives pieces indexed by polymers $Z$ with volume-type decay $e^{-(\kappa'-1)|Z\setminus Y|}$. They depend on $\xi$ and $U'$ only on $Z^{(r')}$, are analytic and invariant, and are bounded in terms of $|\xi|$ on the region where $|Z(s)|_\infty\le\frac12r_I$ and $|\hat Z(Z(s))|_\infty\le\zeta$.
\item \emph{Constants.} The vacuum terms $\mathcal N_k$ and $-\frac12\sum_b\ell_k(b;1)$, and the new main term $g^{-2}S(U')$, are passed on unchanged.
\item \emph{Parameters.} The conditions on the rates are those of Proposition~\ref{prop:localized-pieces} and of Remark~\ref{rem:parameters}, with $\kappa'=\kappa-C\delta$ and the rate $\kappa_n$ of an activity norm. Further, $M>32(L+1)$, $c_a\le c_JL^{-d-2}$, (DC5), and $|\xi|_\infty\le p$ with $gC_Tp\le\frac12r_I$ and $C_wL^dgC_Tp\le\zeta$; these give the two bounds in (a), since $|Z(s)|_\infty\le gC_T|\xi|_\infty$ (Section~\ref{subsec:sq-interaction}) and $|\hat Z(Z)|_\infty\le C_wL^d|Z|_\infty$ (Lemma~\ref{lem:eliminate}(c),(d)). In Section~\ref{sec:cluster} the bound $p$ is the white-noise cutoff $p_0$ and $\zeta=w=C_wL^dC_Tgp_0$; the first condition is part of (Q1), and the second holds with equality. All of these are collected in Remark~\ref{rem:parameters} and Remark~\ref{rem:domain-constants}.
\item \emph{Not done here.} The small-field cutoff and its removal, the Mayer expansion, the bounds in $|\xi|$ and the reblocking to step $k+1$ are in Section~\ref{sec:cluster} and Lemma~\ref{lem:raw-new}; the large-field boundary is not treated.
\end{enumerate}
\end{remark}
